\section*{Chapitre \ref{ch:iv1:waiting-times-and-markov-chains} --- Temps d'attente et chaînes de Markov}
\markboth{\protect\omTocMark{}{Réponses~: Temps d'attente et chaînes de Markov}}{\protect\omTocMark{}{Réponses~: Temps d'attente et chaînes de Markov}}

\begin{solution}{iq:iv1:waiting-times-and-markov-chains:1}
Deux piles consécutives demandent en moyenne 6 lancers. Un pile suivi d'un face en demande 4.

On applique l'analyse au premier pas\index{analyse au premier pas} sur les états de la \cref{fig:iv1:hh-ht}. Pour deux piles, on note $x$ l'espérance du nombre de lancers depuis le départ et $y$ celle du nombre après un pile. Un face depuis l'un ou l'autre état renvoie au départ~:
\[
x = 1 + \tfrac12 x + \tfrac12 y, \qquad y = 1 + \tfrac12 x,
\]
ce qui donne $x = 6$ et $y = 4$. Pour un pile puis un face, on attend d'abord le premier pile, 2 lancers en moyenne, puis le premier face, 2 de plus. Un pile supplémentaire après le premier pile ne fait perdre rien.

\begin{omfigure}
\begin{tikzpicture}[>=stealth,node distance=2.1cm,
  st/.style={circle,draw=omDef,thick,minimum size=11mm,inner sep=0pt,font=\small},
  lab/.style={font=\footnotesize,midway}]
\node[st] (a) {départ};
\node[st,right of=a] (b) {P};
\node[st,right of=b,fill=omDef!15] (c) {PP};
\draw[->] (a) to[bend left=20] node[lab,above] {P} (b);
\draw[->] (b) to[bend left=20] node[lab,below] {F} (a);
\draw[->] (b) -- node[lab,above] {P} (c);
\draw[->] (a) to[out=135,in=225,looseness=5] node[lab,left] {F} (a);
\node[st,right of=c,xshift=1.2cm] (d) {départ};
\node[st,right of=d] (e) {P};
\node[st,right of=e,fill=omDef!15] (f) {PF};
\draw[->] (d) -- node[lab,above] {P} (e);
\draw[->] (e) -- node[lab,above] {F} (f);
\draw[->] (d) to[out=135,in=225,looseness=5] node[lab,left] {F} (d);
\draw[->] (e) to[out=60,in=120,looseness=5] node[lab,above] {P} (e);
\end{tikzpicture}

\omcaption{Attente de PP (à gauche) et de PF (à droite). Un mauvais lancer sur le chemin de PP renvoie au départ, alors qu'un mauvais lancer sur le chemin de PF laisse en place.}\label{fig:iv1:hh-ht}
\end{omfigure}

\iqwrong{«~Chaque motif a une probabilité $1/4$ à chaque position, donc les deux demandent 4 lancers.~» Les deux apparaissent bien une fois toutes les 4 lancers à long terme, mais les occurrences de PP se regroupent en paquets, car PPP en contient deux. Un motif groupé a des intervalles plus longs entre les paquets, donc sa première occurrence prend plus de temps.}
\end{solution}

\begin{solution}{iq:iv1:waiting-times-and-markov-chains:2}
Deux six consécutifs demandent en moyenne 42 lancers. Trois piles consécutifs demandent 14 lancers.

On procède pas à pas avec l'analyse au premier pas\index{analyse au premier pas}. Pour obtenir $k$ à la suite, il faut d'abord obtenir $k - 1$ à la suite, ce qui prend $E_{k-1}$ lancers, puis lancer une fois de plus. Avec la probabilité $p$, c'est gagné~; sinon, on recommence de zéro~:
\[
E_k = E_{k-1} + 1 + (1 - p)\,E_k, \quad\text{donc}\quad E_k = \frac{E_{k-1} + 1}{p}.
\]
Pour le dé, $E_1 = 6$ et $E_2 = (6 + 1) \cdot 6 = 42$. Pour la pièce, $E_1 = 2$, $E_2 = 6$ et $E_3 = 14$.

\iqwrong{36, l'inverse de la probabilité d'une paire de six à une position donnée. C'est la fréquence à long terme, et elle ignore le regroupement qui retarde la première occurrence.}
\iqfollow{Trois six consécutifs~?}{$(42 + 1) \cdot 6 = 258$. La forme générale est $6(6^k - 1)/5$.}
\end{solution}

\begin{solution}{iq:iv1:waiting-times-and-markov-chains:3}
La probabilité de gagner est $6/11 \approx 0.545$.

On applique l'analyse au premier pas\index{analyse au premier pas} sur un tour complet~: on gagne tout de suite avec la probabilité $1/6$, ou bien les deux joueurs ratent, avec la probabilité $(5/6)^2$, et le jeu repart de zéro comme au début, donc
\[
P = \tfrac16 + \bigl(\tfrac56\bigr)^2 P, \quad\text{d'où}\quad P = \frac{1/6}{11/36} = \frac{6}{11}.
\]
\iqfollow{Vous lancez une pièce dont la probabilité de pile est $1/6$, et votre adversaire une pièce dont la probabilité de pile est $1/3$. Vous commencez, et le premier pile gagne. Quelle est votre chance~?}{$\frac{1/6}{1 - (5/6)(2/3)} = 3/8$. Commencer ne compense pas la pièce moins favorable.}
\end{solution}

\enlargethispage{-3\baselineskip}%
\begin{solution}{iq:iv1:waiting-times-and-markov-chains:4}
A gagne avec la probabilité $9/13 \approx 0.692$, et le jeu dure en moyenne $50/13 \approx 3.85$ points de plus.

On regarde le jeu deux points à la fois, avec l'analyse au premier pas\index{analyse au premier pas}. À partir d'un score égal, A gagne les deux avec $p^2 = 0.36$, B gagne les deux avec $q^2 = 0.16$, et sinon les joueurs sont à nouveau à égalité, avec la probabilité $2pq = 0.48$. Donc
\[
W = p^2 + 2pq\,W \;\Rightarrow\; W = \frac{p^2}{p^2 + q^2} = \frac{0.36}{0.52} = \frac{9}{13},
\]
\[
L = 2 + 2pq\,L \;\Rightarrow\; L = \frac{2}{p^2 + q^2} = \frac{50}{13}.
\]
\iqwhy{La règle amplifie un avantage~: un point à 60~\% devient un jeu à 69~\%. C'est l'exemple le plus simple de la façon dont la répétition avec une règle d'arrêt transforme un petit avantage en un grand, au prix d'un jeu plus long.}
\end{solution}

\begin{solution}{iq:iv1:waiting-times-and-markov-chains:5}
Le processus passe un quart de son temps en B, une visite en B dure $10/3$ pas en moyenne, et la corrélation vaut $0.6^n$.

Pour la loi stationnaire\index{loi stationnaire}, le flux de A vers B doit compenser le flux retour~: $\pi_A \cdot 0.1 = \pi_B \cdot 0.3$, donc $\pi_A = 3/4$ et $\pi_B = 1/4$. Une visite en B se termine avec la probabilité $0.3$ à chaque pas, donc sa durée suit une loi géométrique de moyenne $1/0.3 = 10/3$.

Pour la corrélation, une chaîne à deux états a pour valeurs propres 1 et $\lambda = 0.9 + 0.7 - 1 = 0.6$. L'écart à la loi stationnaire se réduit de $\lambda$ à chaque pas, donc la corrélation d'ordre $n$ vaut $\lambda^n = 0.6^n$.

\iqfollow{Un bit traverse 10 copieurs, et chacun le retourne avec la probabilité $0.1$. Quelle est la chance qu'il arrive correct~?}{C'est la même chaîne avec $\lambda = 0.8$~: $\tfrac12(1 + 0.8^{10}) \approx 0.554$.}
\iqwhy{C'est le modèle de régime à deux états. Compte tenu de la persistance de chaque état, il donne la part du temps passée dans l'état, la durée d'un séjour et la décroissance de tout signal fondé sur l'état. Ici, la décroissance de $0.6$ par pas signifie une demi-vie d'un peu plus d'un pas, alors qu'un séjour en A dure 10 pas.}
\end{solution}

\begin{solution}{iq:iv1:waiting-times-and-markov-chains:6}
La probabilité de gagner est $1/4$.

On réfléchit au moment où pile-pile-face peut gagner. Si les deux premiers lancers sont des piles, on gagne à coup sûr~: le premier face qui suit complète PPF, et FPP exige un face plus tôt. Cela a la probabilité $1/4$. Dans tous les autres cas, un face apparaît avant le premier PP. Le premier PP qui apparaît est alors précédé d'un face, et donc FPP se complète en premier.

\iqwrong{Une moitié, par symétrie. Les motifs sont des images miroir l'un de l'autre, mais une course ne lit les lancers que dans le sens direct, donc la symétrie entre eux ne tient pas.}
\iqfollow{Et avec une pièce qui montre pile avec la probabilité $2/3$~?}{Le raisonnement ne change pas~: PPF ne gagne qu'en s'ouvrant par PP, donc sa chance vaut $(2/3)^2 = 4/9$.}
\iqfollow{Combien de temps dure le jeu~?}{$13/2$ lancers en moyenne, moins que les 8 lancers qu'il faut à chacun des deux motifs pris isolément.}
\end{solution}

\begin{solution}{iq:iv1:waiting-times-and-markov-chains:7}
PFPFP demande 42 lancers en moyenne, et PPPFF en demande 32. L'écart vient de la façon dont chaque motif se recouvre avec lui-même.

La formule de Conway\index{formule de Conway} affirme que l'attente est la somme des $2^k$ sur tous les $k$ tels que les $k$ dernières lettres sont égales aux $k$ premières. PFPFP se recouvre avec lui-même pour $k = 1, 3, 5$, ce qui donne $2 + 8 + 32 = 42$. PPPFF ne se recouvre que pour $k = 5$, ce qui donne 32.

Pour le voir, on utilise un argument de casino\index{argument de casino de martingale}. Avant chaque lancer, un nouveau joueur arrive avec 1~\$ et le mise sur le fait que le lancer suivant sera la première lettre du motif. Un gagnant mise les 2~\$ sur la deuxième lettre, et ainsi de suite, jusqu'à ce que le joueur perde ou complète le motif. Tous les paris sont équitables, donc à l'instant d'arrêt, le profit espéré du casino est nul~: les $T$~dollars prélevés aux joueurs égalent en espérance ce qu'il verse. À la fin, les seuls joueurs qui détiennent encore de l'argent sont ceux qui ont commencé il y a $k$ lancers, les $k$ derniers lancers étant les $k$ premières lettres du motif. Chacun détient $2^k$. Donc $E[T] = \sum 2^k$ sur ces $k$.

\iqfollow{Avec un dé, la paire 4 puis 5 contre la paire 4 puis 4~?}{$6^2 = 36$ lancers, contre $36 + 6 = 42$.}
\iqfollow{Un singe tape des majuscules uniformes. Combien de temps faut-il pour ABRACADABRA~?}{Le mot se recouvre avec lui-même pour les longueurs 1, 4 et 11, donc $26^{11} + 26^4 + 26$ frappes.}
\end{solution}

\enlargethispage{-4\baselineskip}%
\begin{solution}{iq:iv1:waiting-times-and-markov-chains:8}
Il faut 10 pas en moyenne.

On utilise la symétrie pour regrouper les coins\index{regroupement des états} selon leur distance à la cible, comme dans la \cref{fig:iv1:cube}. Soient $a$, $b$, $c$ le nombre moyen de pas depuis une distance 3, 2 et 1. Depuis une distance 3, chaque arête mène à une distance 2. Depuis une distance 2, deux des trois arêtes rapprochent de la cible. Depuis une distance 1, une seule des trois mène à la cible~:
\[
a = 1 + b, \qquad b = 1 + \tfrac13 a + \tfrac23 c, \qquad c = 1 + \tfrac23 b.
\]
Ces relations donnent $b = 9$, $c = 7$ et $a = 10$.

\begin{omfigure}
\begin{tikzpicture}[>=stealth,
  st/.style={circle,draw=omDef,thick,minimum size=9mm,inner sep=0pt,font=\small},
  lab/.style={font=\footnotesize,midway}]
\node[st] (a) at (0,0) {3};
\node[st] (b) at (3,0) {2};
\node[st] (c) at (6,0) {1};
\node[st,fill=omDef!15] (d) at (9,0) {0};
\draw[->] (a) to[bend left=20] node[lab,above] {$1$} (b);
\draw[->] (b) to[bend left=20] node[lab,below] {$\tfrac13$} (a);
\draw[->] (b) to[bend left=20] node[lab,above] {$\tfrac23$} (c);
\draw[->] (c) to[bend left=20] node[lab,below] {$\tfrac23$} (b);
\draw[->] (c) -- node[lab,above] {$\tfrac13$} (d);
\node[below=3mm of a,font=\footnotesize] {départ};
\node[below=3mm of d,font=\footnotesize] {cible};
\end{tikzpicture}

\omcaption{La distance de la fourmi au coin opposé. Les huit coins du cube se réduisent à quatre états, contenant respectivement un, trois, trois et un coin.}\label{fig:iv1:cube}
\end{omfigure}

\iqfollow{Combien de temps faut-il pour revenir au départ~?}{8 pas. La marche est uniforme sur les 8 coins à long terme, et un état de part à long terme $\pi$ revient toutes les $1/\pi$ étapes en moyenne.}
\iqfollow{Combien de temps faut-il pour atteindre un coin voisin~?}{7 pas. Le temps d'aller-retour sur une arête vaut $2 \cdot 12 \cdot \tfrac{7}{12} = 14$, et par symétrie la moitié de ce temps va dans chaque sens.}
\iqfollow{Et sur un cube à quatre dimensions~?}{La même chaîne des distances donne $64/3 \approx 21.33$.}
\end{solution}

\begin{solution}{iq:iv1:waiting-times-and-markov-chains:9}
En moyenne, il faut 6 lancers.

Il n'est pas nécessaire d'écrire de système d'équations. Le reste du total modulo 6 est uniforme après chaque lancer, quel que soit le passé, car les six faces atteignent chacun des six restes exactement une fois, depuis n'importe quel reste de départ. Ainsi chaque lancer rend le total divisible par 6 avec probabilité $1/6$, indépendamment du passé, et l'attente suit une loi géométrique de moyenne 6.

\iqfollow{Et divisible par 7~?}{7 lancers. Le premier lancer ne peut jamais réussir, puisqu'il donne de 1 à 6. Ensuite, depuis tout reste non nul, exactement une face atteint un multiple de 7, et un échec laisse sur un autre reste non nul. Il faut donc un lancer, plus une attente géométrique de moyenne 6.}
\iqfollow{Pourquoi cela ne fonctionne-t-il pas pour 5~?}{Cinq restes et six faces ne se répartissent pas également, donc le reste n'est pas uniforme, et il faut résoudre les équations de l'analyse au premier pas\index{analyse au premier pas}.}
\end{solution}

\enlargethispage{-3\baselineskip}%
\begin{solution}{iq:iv1:waiting-times-and-markov-chains:10}
Pour atteindre 10 ou plus, il faut environ $3.32$ lancers. Pour atteindre 100 ou plus, il faut environ $29.05$ lancers.

La méthode exacte est une récurrence par analyse au premier pas\index{analyse au premier pas}. Soit $T(n)$ le nombre moyen de lancers nécessaires pour ajouter au moins $n$. Un lancer donnant $k$ laisse $n - k$ à parcourir, et $T(n) = 0$ pour $n \le 0$~:
\[
T(n) = 1 + \frac16 \sum_{k=1}^{6} T(n - k).
\]
Ainsi $T(1) = 1$, $T(2) = 7/6$, $T(3) = 49/36$, et la table se poursuit jusqu'à $T(10) \approx 3.32$.

Pour vérifier, on utilise l'identité de Wald\index{identité de Wald}~: le total à l'arrêt a pour moyenne $3.5\,E[T]$, et ce total vaut $n$ plus un dépassement. Pour $n$ grand, ce dépassement a pour moyenne
\[
\frac{E[X^2]}{2E[X]} - \frac12 = \frac{91/6}{7} - \frac12 = \frac53.
\]
Donc $E[T] \approx (n + 5/3)/3.5$, ce qui donne $3.33$ pour $n = 10$, et $29.05$ pour $n = 100$.

\iqwrong{$n/3.5$, soit $2.86$ pour $n = 10$. On oublie que l'on s'arrête au-delà de la cible, et non sur celle-ci.}
\end{solution}

\begin{solution}{iq:iv1:waiting-times-and-markov-chains:11}
La probabilité vaut $7/13 \approx 0.538$.

Chaque lancer de la paire donne un 12 avec probabilité $1/36$, un 7 avec probabilité $1/6$, et autre chose avec probabilité $29/36$. La seule mémoire nécessaire est de savoir si la dernière somme était un 7. Soit $a$ la probabilité que 12 sorte en premier lorsque la dernière somme n'était pas un 7, et $b$ celle lorsqu'elle l'était. Par analyse au premier pas\index{analyse au premier pas}, un 7 dans le second état termine la partie par une défaite~:
\[
a = \tfrac{1}{36} + \tfrac16 b + \tfrac{29}{36} a, \qquad b = \tfrac{1}{36} + \tfrac{29}{36} a.
\]
En reportant $b$ dans la première équation et en multipliant par 216, on obtient $42a = 7 + 29a$, donc $a = 7/13$.

\iqwrong{Comparer la probabilité d'obtenir un 12 par lancer, $1/36$, à celle d'obtenir deux 7 de suite, $1/36$, et répondre $1/2$. Deux 7 consécutifs exigent deux lancers qui se chevauchent, ce qui les rend plus lents à arriver qu'un 12.}
\end{solution}

\begin{solution}{iq:iv1:waiting-times-and-markov-chains:12}
L'attente moyenne est de $6.25$ minutes, et non 5.

C'est le paradoxe de l'inspection\index{paradoxe de l'inspection}. Un instant aléatoire tombe plutôt dans un long intervalle~: trois fois plus de temps d'horloge se trouve dans les intervalles de 15 minutes que dans ceux de 5 minutes. On tombe donc dans un intervalle de 15 minutes avec probabilité $3/4$, et l'on attend alors la moitié de celui-ci en moyenne~:
\[
\tfrac34 \cdot 7.5 + \tfrac14 \cdot 2.5 = 6.25.
\]
La formule générale est $E[G^2]/(2E[G])$, soit ici $125/20$.

\iqfollow{Et si les intervalles entre bus suivent une loi exponentielle de moyenne 10 minutes~?}{L'attente vaut 10 minutes, soit la moyenne complète de l'intervalle. L'intervalle dans lequel on tombe a en moyenne 20 minutes. Pour des bus réguliers toutes les 10 minutes, l'attente vaut 5.}
\iqfollow{Dans une longue suite de lancers de pièce, quelle est la longueur de la série qui contient le lancer numéro 1\,000~?}{3 en moyenne, alors qu'une série typique mesure 2. La même pondération par la taille s'applique~: $E[L^2]/E[L] = 6/2$.}
\iqwhy{Toute statistique échantillonnée à des instants aléatoires surreprésente trop souvent les longs intervalles. Mesurer le temps entre les transactions depuis une horloge aléatoire le surestime, et mesurer la fourchette à des instants aléatoires surpondère les carnets d'ordres calmes. Il faut décider si l'on veut une moyenne par événement ou par unité de temps.}
\end{solution}

\begin{solution}{iq:iv1:waiting-times-and-markov-chains:13}
La colonie finit par s'éteindre avec probabilité $1/2$.

On conditionne sur la première étape\index{processus de branchement}. Soit $q$ la probabilité d'extinction. Une famille ne s'éteint que si chaque bactérie vivante après la première minute donne naissance à une famille qui s'éteint, ces familles étant indépendantes~:
\[
q = \tfrac14 + \tfrac14 q + \tfrac12 q^2, \quad\text{ou}\quad 2q^2 - 3q + 1 = 0.
\]
Les racines sont $q = 1/2$ et $q = 1$. Le nombre moyen de descendants vaut $\tfrac14 + 2 \cdot \tfrac12 = 5/4$, supérieur à 1, et dans ce cas la probabilité d'extinction est la plus petite racine. Si la moyenne est au plus égale à 1 (et qu'une bactérie n'est pas simplement remplacée à chaque fois), l'extinction est certaine.

\iqwrong{«~$q = 1$ résout l'équation, donc l'extinction est certaine.~» Cette racine est toujours solution. La probabilité d'extinction est la plus petite racine dans $[0, 1]$, limite des probabilités d'extinction à chaque génération, en partant de zéro.}
\iqfollow{Et en partant de trois bactéries~?}{Les familles sont indépendantes, donc $(1/2)^3 = 1/8$.}
\end{solution}

\begin{solution}{iq:iv1:waiting-times-and-markov-chains:14}
Le total final le plus probable est exactement 1\,000, avec une probabilité d'environ $2/7 \approx 0.286$.

Loin du départ, la probabilité que le total courant atteigne un nombre donné se stabilise à $u = 1/3.5 = 2/7$, par l'argument de renouvellement\index{argument de renouvellement}~: chaque lancer fait avancer le total de $3.5$ en moyenne, donc à long terme un nombre sur $3.5$ est visité.

Pour finir à $1000 + k$, pour $k = 0, \dots, 5$, il faut passer par l'un des nombres de $1000 + k - 6, \dots, 999$, puis obtenir exactement la bonne face. Il y a $6 - k$ tels nombres, chacun visité avec probabilité $2/7$, et chacun exige une face précise, de probabilité $1/6$. Donc
\[
P(\text{total final} = 1000 + k) = (6 - k) \cdot \frac27 \cdot \frac16 = \frac{6 - k}{21}.
\]
Pour $k = 0$, cela vaut $6/21 = 2/7$, la plus grande valeur. Les probabilités valent $6/21, 5/21, \dots, 1/21$, en décroissant jusqu'à la dernière, pour 1\,005. Leur dépassement moyen vaut $5/3$.

\iqwrong{«~Uniforme de 1\,000 à 1\,005, donc $1/6$.~» Cela supposerait que chaque total final soit atteint aussi souvent, alors que 1\,000 peut l'être depuis six nombres, et 1\,005 depuis un seul.}
\iqfollow{Quelle est la probabilité que le lancer qui franchit 1\,000 soit un 2~?}{Par le même décompte, $2 \cdot \tfrac27 \cdot \tfrac16 = 2/21$. Les grands lancers sont plus susceptibles d'être celui qui franchit 1\,000~: la probabilité pour un lancer de $j$ vaut $j/21$.}
\end{solution}

\begin{solution}{iq:iv1:waiting-times-and-markov-chains:15}
Contre HHH, on annonce THH et on gagne avec la probabilité $7/8$. Contre HTH, on annonce HHT et on gagne avec $2/3$. Quel que soit le motif annoncé par l'adversaire, il existe une réponse qui donne au moins $2/3$.

C'est le jeu de Penney\index{jeu de Penney}. La recette consiste à prendre les deux premières lettres du motif de l'adversaire comme les deux dernières du sien, et à commencer par le contraire de la deuxième lettre de l'adversaire. HHH ne gagne que si les trois premiers lancers sont tous des piles, puisqu'une face avant cela fait d'abord apparaître THH. Cela donne à THH la probabilité $7/8$.

En général, on utilise les nombres de tête de Conway\index{formule de Conway}~: $L(X, Y)$ est la somme de $2^{k-1}$ sur tous les $k$ pour lesquels les $k$ dernières lettres de $X$ égalent les $k$ premières de $Y$. Les cotes pour que $B$ batte $A$ valent alors
\[
\bigl(L(A, A) - L(A, B)\bigr) : \bigl(L(B, B) - L(B, A)\bigr).
\]
Pour $A = $ HTH et $B = $ HHT, $L(A, A) = 5$, $L(A, B) = 1$, $L(B, B) = 4$ et $L(B, A) = 2$, donc les cotes valent $4:2$, soit $2/3$.

Le tableau complet des meilleures réponses est~: HHH contre THH ($7/8$), HHT contre THH ($3/4$), HTH contre HHT ($2/3$), HTT contre HHT ($2/3$), et le même avec pile et face échangés.

\iqwrong{«~Choisir un motif qui apparaît aussi souvent, pour que le jeu soit équitable.~» Les huit apparaissent aussi souvent. La relation n'est même pas transitive~: THH bat HHT, qui bat HTT, qui bat TTH, qui bat THH.}
\end{solution}

\begin{solution}{iq:iv1:waiting-times-and-markov-chains:16}
La probabilité vaut $1/11$. La marche entière prend $66$ pas en moyenne.

Tout nombre autre que le point de départ a la même probabilité d'être visité en dernier. Pour que le 6 soit le dernier, la marche doit d'abord atteindre l'un de ses deux voisins, 5 ou 7~; elle le doit, puisqu'elle visite tous les nombres. Supposons qu'elle atteigne d'abord 5. Alors le 6 est le dernier exactement si la marche va de 5 jusqu'à 7 en faisant le tour, avant de poser le pied sur 6. Sur la ligne des positions allant de 6 à 7 en faisant le tour, la marche se trouve à un pas de l'une des extrémités. Une ruine du joueur\index{ruine du joueur} pour une marche équilibrée avec 11 pas entre les extrémités donne la probabilité d'atteindre l'extrémité éloignée~: $1/11$. Rien de cela ne dépend du nombre choisi, donc les 11 ont la même probabilité.

Pour compter les pas, à chaque instant les nombres visités forment un arc. Avec un arc de $j$ nombres, on se trouve à l'une de ses extrémités, et le nouveau nombre suivant prend $j$ pas en moyenne~: c'est une marche équilibrée sur $j + 2$ positions qui part d'à côté de l'une des extrémités, et une telle marche prend $1 \cdot j$ pas. En sommant de $j = 1$ à 11 on obtient $66$, soit $n(n-1)/2$ en général.

\iqwrong{Supposer que le 6, étant le plus éloigné, soit le plus susceptible d'être le dernier. Il met le plus de temps à l'atteindre en moyenne, soit $6 \cdot 6 = 36$ pas, mais c'est une autre question.}
\end{solution}

\begin{solution}{iq:iv1:waiting-times-and-markov-chains:17}
Il faut 168 coups en moyenne.

Aucun calcul sur l'échiquier n'est nécessaire. Pour une marche aléatoire sur un graphe, qui choisit à chaque pas une arête uniformément au hasard, la loi stationnaire\index{loi stationnaire} est proportionnelle au degré, c'est-à-dire au nombre d'arêtes issues d'un sommet~:
\[
\pi(v) = \frac{\deg(v)}{\sum_u \deg(u)}.
\]
On le vérifie en une ligne~: le flux entrant en $v$ vaut $\sum_{u \sim v} \pi(u)/\deg(u)$, qui comporte $\deg(v)$ termes égaux, chacun égal à $1/\sum_u \deg(u)$. Le temps de retour moyen dans un état vaut $1/\pi(v)$, d'après le lemme de Kac\index{lemme de Kac}.

Comptons maintenant les coups du cavalier. Un coup de cavalier a 2 orientations et 4 choix de signes, et chacun des 8 types a $7 \cdot 6$ placements sur l'échiquier, donc $\sum_u \deg(u) = 8 \cdot 42 = 336$. Une case d'angle a un degré 2, donc le retour prend $336/2 = 168$ coups.

\iqwrong{Essayer de résoudre 63 équations linéaires sur l'échiquier. Un recruteur qui pose cette question attend l'argument sur les degrés.}
\iqfollow{Et depuis une case centrale~?}{Degré 8, donc $336/8 = 42$ coups.}
\iqfollow{Et pour un roi partant d'une case d'angle~?}{La somme des degrés du roi vaut 420, et une case d'angle a un degré 3, donc 140 coups.}
\end{solution}

\begin{solution}{iq:iv1:waiting-times-and-markov-chains:18}
Il faut 18 pas en moyenne.

La réflexion par rapport à la droite joignant le départ à la cible regroupe les positions\index{regroupement des états} en classes. L'analyse au premier pas\index{analyse au premier pas} sur ces classes donne les valeurs dans la \cref{fig:iv1:grid}~:
\[
s = 1 + a, \quad a = 1 + \tfrac13(s + m + c), \quad m = 1 + \tfrac14(2a + 2b), \quad c = 1 + \tfrac12(a + b), \quad b = 1 + \tfrac13(m + c),
\]
où $s$, $a$, $m$, $c$, $b$ désignent le départ, ses voisins, le centre, les deux autres coins et les voisins de la cible. Les équations pour $m$ et $c$ ont le même second membre, donc $m = c$. La solution est $s = 18$, $a = 17$, $m = c = 15$ et $b = 11$.

\begin{omfigure}
\begin{tikzpicture}[x=1.7cm,y=1.7cm,
  st/.style={circle,draw=omDef,thick,minimum size=8mm,inner sep=0pt,font=\small,fill=white}]
\foreach \i in {0,1,2}{\draw[omDef] (\i,0) -- (\i,2); \draw[omDef] (0,\i) -- (2,\i);}
\node[st] at (0,0) {18};
\node[st] at (1,0) {17};
\node[st] at (2,0) {15};
\node[st] at (0,1) {17};
\node[st] at (1,1) {15};
\node[st] at (2,1) {11};
\node[st] at (0,2) {15};
\node[st] at (1,2) {11};
\node[st,fill=omDef!15] at (2,2) {0};
\node[below=5mm,font=\footnotesize] at (0,0) {départ};
\node[above=5mm,font=\footnotesize] at (2,2) {cible};
\end{tikzpicture}

\omcaption{Nombre moyen de pas pour atteindre le coin supérieur droit depuis chaque position de la grille. Chaque valeur vaut un de plus que la moyenne sur ses voisins.}\label{fig:iv1:grid}
\end{omfigure}

Une méthode plus élégante est l'identité du temps d'aller-retour\index{identité du temps d'aller-retour}~: on traite chaque liaison comme une résistance de 1~ohm. Le temps moyen d'un aller-retour entre deux sommets vaut alors $2 \cdot (\text{nombre d'arêtes}) \cdot R$, où $R$ est la résistance effective. Ici $R = 3/2$, donc un aller-retour prend $2 \cdot 12 \cdot 3/2 = 36$ pas. Par symétrie, chaque trajet prend la moitié de ce temps.

\iqfollow{Le temps est-il toujours le même dans chaque sens~?}{Non. Du centre vers un coin, il faut 15 pas, et d'un coin vers le centre seulement 6. Leur somme, 21, vaut $2 \cdot 12 \cdot 7/8$, comme le dit l'identité.}
\end{solution}

\begin{solution}{iq:iv1:waiting-times-and-markov-chains:19}
La variance vaut 22, donc l'écart-type vaut environ $4.7$, proche de la moyenne 6. Pour pile puis face, la moyenne vaut 4 et la variance n'est que de 4.

On applique l'analyse au premier pas\index{analyse au premier pas} au moment d'ordre 2\index{moment d'ordre 2}. Avec les moyennes $m_0 = 6$ depuis le départ et $m_1 = 4$ après un pile, posons $a = E[T_0^2]$ et $b = E[T_1^2]$. Un lancer ajoute un au temps restant, et $E[(1 + T)^2] = 1 + 2E[T] + E[T^2]$. Donc
\[
a = \tfrac12(1 + 2 \cdot 6 + a) + \tfrac12(1 + 2 \cdot 4 + b) = 11 + \tfrac12 a + \tfrac12 b,
\]
\[
b = \tfrac12 \cdot 1 + \tfrac12(1 + 2 \cdot 6 + a) = 7 + \tfrac12 a.
\]
On en tire $a = 58$ et $b = 36$, donc la variance vaut $58 - 6^2 = 22$.

L'attente de pile puis face est la somme de deux attentes géométriques indépendantes, de moyenne 2 et de variance 2, donc sa variance vaut 4.

\iqfollow{Et pour trois piles d'affilée~?}{La même méthode donne une moyenne de 14 et une variance de 142.}
\iqwhy{Deux attentes de moyennes voisines peuvent présenter des risques très différents. Une attente dont les échecs ramènent au début a un écart-type proche de sa moyenne, donc il faut prévoir une longue queue de distribution et ne pas citer la moyenne seule.}
\end{solution}

\begin{solution}{iq:iv1:waiting-times-and-markov-chains:20}
Il faut $52\,H_{52} \approx 236$ mouvements en moyenne, où $H_n = \sum_{k \le n} 1/k$.

Suivons la carte qui était au fond du paquet au départ. Lorsque $i$ cartes se trouvent sous elle, la carte du dessus vient s'intercaler sous elle avec probabilité $(i + 1)/52$, car il y a $i + 1$ positions sous elle. L'attente moyenne pour que cette carte remonte d'une place vaut donc $52/(i + 1)$. Elle atteint le dessus lorsque $i = 51$, et un dernier mouvement l'insère~:
\[
\sum_{i=0}^{50} \frac{52}{i+1} + 1 = 52\,H_{51} + 1 = 52\,H_{52} \approx 236.
\]
C'est le problème du collectionneur de coupons\index{collectionneur de coupons} sous un autre visage.

Le jeu est alors parfaitement mélangé, par récurrence\index{temps uniforme fort}. Les cartes situées sous la carte marquée sont toujours dans un ordre uniformément aléatoire~: une nouvelle carte se place uniformément à l'une des $i + 1$ places parmi les $i$ cartes déjà présentes. Lorsque la carte marquée atteint le dessus, les 51 cartes qui se trouvent dessous sont dans un ordre uniforme. La dernière insertion la place uniformément parmi les 52 positions, donc tout le jeu est uniforme.

\iqfollow{236 est-il le temps de mélange~?}{C'est une borne supérieure en espérance~: le jeu peut déjà être proche de l'uniforme plus tôt. Elle indique l'ordre de grandeur du nombre de mouvements à prévoir~: $n \ln n$ mouvements.}
\end{solution}

\begin{solution}{iq:iv1:waiting-times-and-markov-chains:21}
Atteindre l'état «~que des piles~» prend $74752/63 \approx 1186.5$ pas en moyenne. Revenir à l'état «~que des faces~» prend exactement 1\,024.

Le temps de retour ne demande presque aucun calcul. L'état est un sommet d'un cube de dimension 10, et la marche admet une loi stationnaire\index{loi stationnaire} uniforme sur les $2^{10}$ sommets, puisque chaque sommet a le même degré. Par le lemme de Kac\index{lemme de Kac}, le temps de retour moyen vaut donc $2^{10} = 1024$.

Pour atteindre «~que des piles~», on ne suit que le nombre de faces. Soit $e_k$ l'espérance du nombre de coups pour passer de $k$ faces à $k - 1$. Avec la probabilité $k/10$, on obtient une face et c'est fini. Sinon, on crée une face et il faut redescendre de $k + 1$, puis de nouveau de $k$~:
\[
e_k = 1 + \frac{10 - k}{10}(e_{k+1} + e_k) \quad\Rightarrow\quad e_k = \frac{10}{k} + \frac{10 - k}{k}\,e_{k+1},
\]
avec $e_{10} = 1$. En descendant et en sommant $e_{10} + \dots + e_1$, on obtient $74752/63$. La même récurrence avec 3 pièces donne 10, la fourmi du cube.

\iqwrong{Supposer que l'état le plus éloigné demande un temps beaucoup plus long. Les deux temps sont d'ordre $2^{10}$, car l'essentiel du temps est passé à errer autour du milieu, vers cinq faces, avant une rare série jusqu'à l'un des extrêmes.}
\end{solution}

\begin{solution}{iq:iv1:waiting-times-and-markov-chains:22}
Nommons THTH. On gagne avec la probabilité $9/14 \approx 0.643$, et le jeu dure $90/7 \approx 12.86$ lancers en moyenne.

Utilisons les nombres de tête de Conway\index{formule de Conway} avec deux casinos\index{argument de casino par martingale}. Comme pour un seul motif, un joueur qui a commencé il y a $k$ lancers détient $2^k$ si les $k$ derniers lancers coïncident avec les $k$ premières lettres du motif sur lequel il mise. On fait tourner un casino de joueurs qui misent sur $A$ et un autre qui misent sur $B$. Soit $L(X, Y) = \sum 2^{k-1}$ sur les $k$ pour lesquels les $k$ dernières lettres de $X$ sont les $k$ premières de $Y$, et soit $p$ la probabilité que $A$ gagne. Chaque casino est équitable en espérance, donc
\[
E[T] = 2p\,L(A,A) + 2(1-p)\,L(B,A) = 2p\,L(A,B) + 2(1-p)\,L(B,B).
\]
En égalant les deux membres de droite, on obtient la cote de $B$ contre $A$, $\;L(A,A) - L(A,B) : L(B,B) - L(B,A)$.

Les meilleures réponses commencent par une lettre, puis recopient les trois premières lettres de $A$, ici $x$HTH, de sorte que $B$ ait besoin d'un lancer de moins chaque fois que $A$ est en cours. Pour $A = $ HTHH et $B = $ THTH, $L(A,A) = 9$, $L(A,B) = 0$, $L(B,B) = 10$ et $L(B,A) = 5$. La cote est $9:5$, soit $9/14$. Puis $E[T] = 2 \cdot \tfrac{5}{14} \cdot 9 + 2 \cdot \tfrac{9}{14} \cdot 5 = 90/7$. Le second casino donne le même résultat~: $2 \cdot \tfrac{9}{14} \cdot 10 = 90/7$. L'autre candidat, HHTH, n'obtient que $(9-3):(9-5)$, soit $3/5$.

\iqwrong{Appliquer la recette des trois lancers sans vérification~: commencer par le contraire de la deuxième lettre de $A$, puis recopier les trois premières lettres de $A$. Cela donne HHTH, qui n'est pas le meilleur choix. Pour quatre lancers, la première lettre doit être choisie en calculant les deux possibilités.}
\end{solution}

\begin{solution}{iq:iv1:waiting-times-and-markov-chains:23}
Le nombre espéré de lancers est $416/5 = 83.2$.

Trouvons les bons états\index{regroupement d'états}. Soit l'état $j$ le nombre des derniers lancers qui sont tous différents, de sorte qu'on part de l'état 1 après un lancer et qu'on s'arrête à l'état 6. Depuis l'état $j$, une nouvelle face a la probabilité $(6 - j)/6$ et fait passer à $j + 1$. Sinon, le lancer répète la face apparue il y a $i$ lancers, pour un $i = 1, \dots, j$, chacun avec la probabilité $1/6$, et la nouvelle suite distincte a une longueur exactement égale à $i$, comme dans la \cref{fig:iv1:distinct}.

Soit $e_j$ le temps moyen pour atteindre $j + 1$ pour la première fois depuis $j$. Une chute vers l'état $i$ coûte la remontée, $e_i + \dots + e_j$. En développant le premier pas et en regroupant les termes,
\[
e_j = 1 + \frac16 \sum_{i=1}^{j} \sum_{m=i}^{j} e_m = 1 + \frac16 \sum_{m=1}^{j} m\,e_m \quad\Rightarrow\quad (6 - j)\,e_j = 6 + \sum_{m < j} m\,e_m.
\]
Cela se résout de bas en haut~: $e_1 = \tfrac65$, $e_2 = \tfrac95$, $e_3 = \tfrac{18}{5}$, $e_4 = \tfrac{54}{5}$, $e_5 = \tfrac{324}{5}$. En ajoutant le premier lancer, le total vaut $1 + 411/5 = 416/5$.

\begin{omfigure}
\begin{tikzpicture}[>=stealth,x=1.8cm,
  st/.style={circle,draw=omDef,thick,minimum size=8mm,inner sep=0pt,font=\small},
  lab/.style={font=\footnotesize}]
\foreach \j in {1,...,5}{\node[st] (s\j) at (\j,0) {\j};}
\node[st,fill=omDef!15] (s6) at (6,0) {6};
\draw[->] (s1) -- (s2);
\draw[->] (s2) -- (s3);
\draw[->] (s3) -- (s4);
\draw[->] (s4) -- node[above,lab] {$\tfrac26$} (s5);
\draw[->] (s5) -- (s6);
\draw[->,omThm] (s4) to[bend left=45] node[below,lab,pos=0.5] {$\tfrac16$} (s1);
\draw[->,omThm] (s4) to[bend left=50] node[above,lab,pos=0.5] {$\tfrac16$} (s2);
\draw[->,omThm] (s4) to[bend left=60] node[above,lab,pos=0.5] {$\tfrac16$} (s3);
\draw[->,omThm] (s4) to[out=60,in=120,looseness=5] node[above,lab] {$\tfrac16$} (s4);
\end{tikzpicture}

\omcaption{La chaîne sur le nombre de lancers distincts récents, ne montrant que les transitions depuis l'état 4. Une répétition de la face vue il y a $i$ lancers ramène à l'état $i$. Chaque état se comporte de la même manière.}\label{fig:iv1:distinct}
\end{omfigure}

\iqwrong{$6^6/6! = 64.8$, l'inverse de la probabilité qu'une fenêtre donnée de six lancers soit entièrement distincte. C'est la fréquence à long terme de ces fenêtres, et comme les fenêtres se recouvrent et se regroupent, la première arrive plus tard.}
\end{solution}

\begin{solution}{iq:iv1:waiting-times-and-markov-chains:24}
La famille s'éteint avec la probabilité $1/2$, et sachant cela, le total espéré vaut 3 cellules.

La probabilité d'extinction vérifie $q = \tfrac13 + \tfrac23 q^2$, en conditionnant sur le premier pas\index{processus de branchement}. Les racines sont $1/2$ et 1, et comme la moyenne des descendants, $4/3$, est supérieure à 1, la réponse est la plus petite racine.

Sans conditionnement, l'espérance du total est infinie, car la famille croît sans limite la moitié du temps. L'idée consiste à conditionner par l'extinction\index{conditionnement par l'extinction}. Une cellule qui a $k$ enfants a une famille qui s'éteint avec la probabilité $q^k$, indépendamment pour chaque enfant, de sorte que la loi conditionnelle du nombre d'enfants vaut $P(k \text{ enfants} \mid \text{extinction}) = p_k\,q^k / q$~:
\[
P(0 \mid \text{extinct}) = \frac{1/3}{1/2} = \frac23, \qquad P(2 \mid \text{extinct}) = \frac{(2/3)(1/4)}{1/2} = \frac13.
\]
Sachant qu'elle s'éteint, la famille de chaque enfant est à nouveau une copie du même processus conditionné. Donc, conditionnellement à l'extinction, on a un processus de branchement de moyenne $\tfrac23$. La génération $n$ a une taille espérée de $(2/3)^n$, et le total vaut
\[
\sum_{n \ge 0} \Bigl(\frac23\Bigr)^n = 3.
\]
\iqwrong{Affirmer que le total est infini, ou calculer la moyenne avec les probabilités initiales. Le conditionnement modifie la loi des descendants, et c'est cette étape qui rend la question difficile.}
\iqfollow{Sachant l'extinction, quelle est la probabilité que la première cellule meure immédiatement~?}{$2/3$, la première probabilité ci-dessus, en hausse par rapport à $1/3$.}
\end{solution}

\begin{solution}{iq:iv1:waiting-times-and-markov-chains:25}
La génération 10 est vide avec la probabilité $10/11$. Si elle n'est pas vide, sa taille espérée vaut 11. Le temps d'extinction espéré est infini, bien que l'extinction soit certaine.

\begin{sloppypar}Utilisons la fonction génératrice des probabilités\index{fonction génératrice des probabilités}. La loi des descendants donne $f(s) = \sum_k s^k/2^{k+1} = 1/(2 - s)$, et la génération $n$ a pour fonction génératrice la $n$-ième composée $f_n$ de $f$. L'idée qui supprime les calculs est que cette application devient linéaire après un changement de variable\index{itération linéaire-fractionnaire}~:
\[
\frac{1}{1 - f(s)} = \frac{2 - s}{1 - s} = 1 + \frac{1}{1 - s}.
\]
Chaque itération ajoute donc 1, et $\dfrac{1}{1 - f_n(s)} = n + \dfrac{1}{1 - s}$. En $s = 0$, cela donne $1/(1 - f_n(0)) = n + 1$, donc $P(\text{la génération } n \text{ est vide}) = f_n(0) = n/(n+1)$. Pour $n = 10$, cela vaut $10/11$.\end{sloppypar}

La moyenne des descendants vaut 1, donc $E[X_n] = 1$ pour tout $n$. Sachant que la famille survit, $E[X_n \mid X_n > 0] = 1/P(X_n > 0) = n + 1$, soit 11.

La famille dépasse la génération $n$ avec la probabilité $1/(n + 1)$. La somme de ces queues pour le temps d'extinction espéré donne la série harmonique, qui diverge.

\iqwrong{«~Une moyenne de un descendant signifie une population stable.~» Une famille critique s'éteint avec certitude. La moyenne reste à 1 seulement parce que les rares survivants croissent~: à la génération $n$, une seule famille sur $n + 1$ survit, avec une taille moyenne de $n + 1$.}
\iqwhy{Tout ce qui se reproduit exactement au taux de remplacement se comporte ainsi~: la plupart des lignées disparaissent tôt, et les rares survivantes paraissent grandes. Ne regarder que les survivants surestime la croissance, d'un facteur qui augmente avec le temps.}
\end{solution}
